\begin{table}[t]
\centering
\caption{Standardized Effect Sizes}
\label{tab:sde}
\begin{threeparttable}
\begin{tabular}{lcccccc}
\hline\hline
Outcome & $\hat{\beta}$ & SE & SD($Y$) & SDE & SE(SDE) & Classification \\
\hline
$\Delta$RN 2017--2012 & -0.1660 & 0.0140 & 6.80 & -0.218 & 0.018 & Large neg. \\
$\Delta$RN 2022--2012 & 0.0855 & 0.0083 & 6.08 & 0.126 & 0.012 & Mod. pos. \\
$\Delta$Turnout 2017--2012 & 0.0009 & 0.0037 & 3.80 & 0.002 & 0.009 & Null \\
$\Delta$M\'{e}lenchon 2017--2012 & 0.0050 & 0.0085 & 5.29 & 0.009 & 0.014 & Small pos. \\
\hline\hline
\end{tabular}
\begin{tablenotes}[flushleft]\small
\item \textit{Notes:} \textbf{Country:} France. \textbf{Research question:} Does the 2014--2018 carbon tax escalation (TICPE composante carbone) cause differential increases in National Rally (RN) vote share in car-dependent communes? \textbf{Policy mechanism:} France's carbon component of the domestic fuel tax (TICPE) rose from 7 to 44.60 EUR/tCO\textsubscript{2} over 2014--2018, adding approximately 12 cents per liter to diesel; communes with higher car-commuting shares face a larger effective fuel-cost burden per household. \textbf{Outcome definition:} Change in first-round presidential RN vote share (percentage points), 2017 minus 2012 (primary) and 2022 minus 2012 (long-run). \textbf{Treatment:} Continuous; commune-level car-commuting share from the 2011 census (fraction of employed residents age 15+ commuting by car). \textbf{Data:} Minist\`{e}re de l'Int\'{e}rieur election results (data.gouv.fr), INSEE RP 2011 census (transport mode), Filosofi 2019 (income); 33,390 communes, 2007--2022. \textbf{Method:} Cross-sectional first-difference with d\'{e}partement fixed effects and controls; standard errors clustered at d\'{e}partement level (96 clusters). \textbf{Sample:} Metropolitan France communes with $\geq$50 registered voters and $\geq$10 employed residents in 2011. SDE $= \hat{\beta} \times \text{SD}(X) / \text{SD}(Y)$ where SD($X$) is the cross-commune standard deviation of car-commuting share and SD($Y$) is the outcome standard deviation. Classification refers to magnitude, not statistical significance: Large ($|$SDE$| > 0.15$), Moderate ($0.05$--$0.15$), Small ($0.005$--$0.05$), Null ($< 0.005$).
\end{tablenotes}
\end{threeparttable}
\end{table}

